% docs/manual/user/chapters/06-running.tex
% 第 6 章：运行评估

\chapter{运行评估}

第 2 章的 Quick Start 把 \cli{openbench check} 与 \cli{openbench run} 当成黑盒。本章打开盒子：每个命令具体在做什么、阶段如何衔接、缓存如何工作、出问题怎么诊断。本章是用户从"会跑"到"能调"的桥梁。

\section{\cli{openbench check} 详解}

\cli{check} 是 \emph{零数据访问} 的预飞行检查，只读 catalog 与 schema：

\begin{minted}{bash}
openbench check openbench.yaml
\end{minted}

执行步骤：

\begin{enumerate}
  \item \textbf{文件存在性}：检查 path 存在且后缀是 \texttt{.yaml/.yml}
  \item \textbf{YAML 语法}：用 PyYAML safe loader 解析（含 \texttt{!include} 展开）
  \item \textbf{Schema 校验}：必填字段齐、字段类型对、枚举值在合法集
  \item \textbf{时段合理性}：\yamlkey{years[0] <= years[1]}，且都是整数
  \item \textbf{Reference 解析}：每个 \yamlkey{evaluation.variables} 在 \yamlkey{reference} 有映射；映射的数据集名能在 registry 找到（含 \texttt{\_LowRes}/\texttt{\_MidRes} 自动解析）
  \item \textbf{Provenance 标注}：每个解析后的 reference，按元数据可信度打 ✓/⚠/\textasciitilde
  \item \textbf{Simulation 列出}：报告每个 simulation 的 \yamlkey{model} 与 \yamlkey{root_dir}（不读 NC 文件，只显示路径）
\end{enumerate}

\subsection{Provenance 标记含义}

每个 reference 在解析时会附带元数据来源标识：

\begin{itemize}
  \item \texttt{✓}（绿）：高置信。值来自数据集 catalog 显式声明或 model profile
  \item \texttt{\textasciitilde}（青）：中置信。从目录结构推断（如 \texttt{Grid/LowRes/} 推出 0.5°）
  \item \texttt{⚠}（黄）：低置信。使用了默认值，可能与实际不符
\end{itemize}

\subsection{strict\_reference 模式}

\begin{minted}{yaml}
project:
  strict_reference: true
\end{minted}

启用后：低置信度 (⚠) 警告升级为错误，\cli{check} 返回非零退出码。适合"必须保证元数据精准"的生产环境。

\section{\cli{openbench run} 阶段流程}

\cli{run} 把已通过 check 的 config 翻译成实际计算。内部分 5 个阶段：

\begin{enumerate}
  \item \textbf{Build}：用 \modname{config.adapter} 把新 schema 翻译成 evaluator 期待的 legacy namelist；准备输出目录树
  \item \textbf{Evaluation}：每个 (variable, simulation, reference) 三元组独立计算 metrics + scores
  \item \textbf{Comparison}：（可选）跨 simulation 合并指标，生成 portrait/radar/ranking 图
  \item \textbf{Statistics}：（可选）跑 ANOVA、相关性、KDE、Mann-Kendall 等
  \item \textbf{Report}：（可选）汇总 HTML/PDF 报告
\end{enumerate}

\subsection{Build 阶段}

把新 \file{openbench.yaml} 翻译成内部 namelist 格式（main\_nl / ref\_nml / sim\_nml / fig\_nml）。\textbf{用户视角}通常无感。

如要看翻译结果（调试用）：

\begin{minted}{bash}
openbench run openbench.yaml --dump-config
# 中间产物会写到 output/<run>/debug/{runner_config,main_nl,ref_nml,sim_nml,fig_nml}.yaml
\end{minted}

\subsection{Evaluation 阶段}

对每个 task（一个 var × 一个 sim × 一个 ref）：

\begin{enumerate}
  \item 加载 reference NC（按 catalog \yamlkey{root_dir}）
  \item 加载 simulation NC（按 \yamlkey{simulation.<label>.root_dir}）
  \item 应用单位换算（profile 中的 \yamlkey{compute} 表达式）
  \item 重网格化到 \yamlkey{project.grid_res}
  \item 时间对齐到 \yamlkey{project.tim_res}（按 \yamlkey{time_alignment} 策略）
  \item 应用 \yamlkey{unified_mask}（统一 NaN 掩码，跨 sim 公平比较）
  \item 计算所有 metric 与 score
  \item 写 CSV + 调用 visualization 模块出图
\end{enumerate}

输出位置：

\begin{itemize}
  \item Metrics：\file{output/<run>/metrics/<var>/<sim>\_metrics.csv}
  \item Scores：\file{output/<run>/scores/<var>/<sim>\_scores.csv}
  \item 图：\file{output/<run>/figures/<var>/<fig\_kind>/<sim>\_<ref>.png}
\end{itemize}

\subsection{Comparison 阶段（条件触发）}

仅当 \yamlkey{comparison.enabled: true} 且至少有 1 个 simulation 时执行：

\begin{itemize}
  \item 把所有 simulation 的 metrics/scores 合并为对比表
  \item 生成 portrait plot（多模型 × 多变量的统一打分网格）
  \item radar plot（每模型一组指标的雷达图）
  \item ranking 表与图
\end{itemize}

输出位置：\file{output/<run>/comparisons/...}

\subsection{Statistics 阶段（条件触发）}

仅当 \yamlkey{statistics.enabled: true} 时执行。跑各类统计模块：

\begin{itemize}
  \item ANOVA、相关、协方差
  \item Mann-Kendall 趋势检验
  \item Hellinger Distance、KDE
  \item 三角帽、偏最小二乘回归
  \item 滑动窗口 / 重采样
\end{itemize}

输出位置：\file{output/<run>/statistics/...}

\subsection{Report 阶段（条件触发）}

仅当 \yamlkey{project.generate_report: true}（默认 \texttt{true}）且安装了 \cli{openbench[report]} 时执行：

\begin{itemize}
  \item 收集所有 figure / metric 表，按章节组织
  \item 渲染 HTML（一定）
  \item 渲染 PDF（如果 \modname{xhtml2pdf} 可用）
\end{itemize}

输出位置：\file{output/<run>/report.html} 与 \file{output/<run>/report.pdf}。

\warninline{未装 \cli{[report]} extra 时，run 仍会成功，但 \file{report.html/pdf} 不会生成。日志会有"report generation skipped"提示。}

\section{常用 CLI 选项}

\begin{minted}{bash}
# 仅校验 + 显示将要做什么；不实际计算
openbench run openbench.yaml --dry-run

# 限制只跑某变量子集（覆盖 evaluation.variables）
openbench run openbench.yaml --variables Evapotranspiration --variables Latent_Heat

# 覆盖并行核数（覆盖 project.num_cores）
openbench run openbench.yaml --cores 16

# 跳过 evaluation，仅基于已有结果重做 comparison
openbench run openbench.yaml --comparison-only

# 输出中间 namelist 与 runner config 到 debug/
openbench run openbench.yaml --dump-config
\end{minted}

\warninline{\cli{--remote HOST} 标志在当前 v3.0a1 \emph{尚未在 CLI 实现}。传入后会立刻退出并提示用 GUI。运维卷第 3 章会详述远程执行的现状与替代方案。}

\section{增量缓存与 force 重算}

\cli{run} 默认启用增量缓存：

\subsection{缓存机制}

每个 task 完成后，\openbench{} 把：

\begin{itemize}
  \item 任务配置（var/sim/ref/metrics/scores 列表）的哈希
  \item 输出文件清单
\end{itemize}

写入 \file{output/<run>/scratch/cache/<task\_key>.json}。

下一次 run 同一 config：

\begin{itemize}
  \item 计算每个 task 的 \emph{当前} 配置哈希
  \item 与 cache 中的哈希对比，相同 → 跳过该 task（直接复用现有输出）
  \item 任一字段变化 → 重新计算
\end{itemize}

实测时长上，增量重跑通常只用首次的 5-15\%。

\subsection{触发全量重算}

\begin{minted}{yaml}
project:
  force: true   # 每次 run 都全量重算
\end{minted}

\textbf{重要}：\openbench{} v3.0 \cli{run} 子命令\emph{不}提供 \texttt{--force} 命令行 flag——只能通过配置文件设置。\modname{runner.local:409} 在没有显式 force 参数传入时回退读取 \yamlkey{cfg.project.force}。临时调试可用环境变量包装或临时改 yaml 后再改回。

\paragraph{缓存哈希覆盖范围}

\modname{runner.cache.EvaluationCache.hash_config} 把以下字段进哈希（\modname{runner.local:421--429}）：

\begin{itemize}
  \item \texttt{variable} / \texttt{sim\_source} / \texttt{ref\_source}
  \item \texttt{metrics} 列表 / \texttt{scores} 列表
  \item \texttt{comparisons} 列表 / \texttt{statistics} 列表
\end{itemize}

\emph{不进}哈希的：\texttt{--cores} / \texttt{--variables} CLI override、\texttt{num\_cores}、\texttt{generate\_report}、\texttt{debug\_mode}。这意味着：

\begin{itemize}
  \item 改 metric/score 列表 → 自动重算受影响的 task（即使没设 \yamlkey{force}）
  \item 改并行核数 → 直接复用缓存，不重算
  \item 改输出目录路径（\yamlkey{output\_dir}）→ 哈希不变但目录变，缓存路径找不到 → 实际重算
\end{itemize}

最直接的全量重算手段：删除 \file{output/<run>/} 整个目录。

\subsection{何时该全量重算}

\begin{itemize}
  \item 升级 \openbench{}（metric/score 实现改变，但版本号没进哈希）
  \item 改了 reference catalog（数据集本身的元数据）
  \item 怀疑缓存被污染（异常中断、文件系统故障后）
\end{itemize}

\section{only\_drawing 模式：只重出图}

如果只想调整可视化（颜色、范围、标题）而保留已有计算结果：

\begin{minted}{yaml}
project:
  only_drawing: true
\end{minted}

CLI flag 未暴露——只能通过 yaml 切换。

\paragraph{实现}

打开后 \openbench{} 切换到 \modname{visualization.Mod\_Only\_Drawing} 下的 5 个专用类：

\begin{itemize}
  \item \modname{Evaluation\_grid\_only\_drawing} —— 替换 \modname{core.evaluation.Evaluation\_grid}
  \item \modname{Evaluation\_stn\_only\_drawing} —— 替换 \modname{Evaluation\_stn}
  \item \modname{LC\_groupby\_only\_drawing} —— 替换 \modname{core.landcover\_groupby.LC\_groupby}
  \item \modname{CZ\_groupby\_only\_drawing} —— 替换 \modname{core.climatezone\_groupby.CZ\_groupby}
  \item \modname{ComparisonProcessing\_only\_drawing} —— 替换 \modname{core.comparison.ComparisonProcessing}
\end{itemize}

每个类与原版同接口，但跳过计算，直接从磁盘读 \file{metrics/}、\file{scores/}、\file{comparisons/} 下已有的 NC/CSV，调 visualization 模块重出图。

\paragraph{文件依赖}

\begin{itemize}
  \item \file{output/<run>/data/<var>\_ref\_<R>\_<rv>.nc} —— 预处理后的 ref（mass weight 用）
  \item \file{output/<run>/metrics/}、\file{scores/}、\file{comparisons/} 完整的逐 task NC
  \item \file{output/<run>/scratch/cache/} 仅 cache 状态参考，不影响出图
\end{itemize}

任一缺失会打 \texttt{ERROR}：\emph{"File not found in only\_drawing mode: ... Please run the evaluation first (set only\_drawing=False) to generate required data files."} 并跳过该图。常见触发：先开 \yamlkey{only\_drawing} 看效果但忘了先跑过完整评估。

\section{debug\_mode 与日志}

\subsection{开启 debug}

\begin{minted}{yaml}
project:
  debug_mode: true
\end{minted}

\textbf{准确范围}（核对自代码）：\yamlkey{debug_mode} \emph{不}是全局日志级别开关。它只影响 \modname{data.processing.DatasetProcessing} 与若干 \modname{data.custom.*\_filter}（如 HydroWeb）内部的诊断打印：

\begin{itemize}
  \item \modname{processing.py:730}：缺失关键属性时打 WARNING 提示用了哪个 fallback（\yamlkey{tim_res} 替代 \yamlkey{compare_tim_res} 等）
  \item \modname{processing.py:2402}：站点处理某个分支的额外诊断
  \item HydroWeb / GEBA 等 custom filter：站点筛选时打印每个站的过滤理由
\end{itemize}

不影响：runner 阶段日志（\modname{runner.local}）、preprocess 进度、unified\_mask、cache hit/miss 这些都是 \texttt{logger.info} / \texttt{logger.warning} 直接走 \openbench{} 的根 logger，与 \yamlkey{debug_mode} 无关。

要看运行期完整 debug 输出，应该用 Python 标准 logging 配置（开发卷第 4 章会给具体方法），或者在 CLI 加 \texttt{2>\&1 | grep DEBUG}（但这只在你已用 logging.basicConfig 设过 DEBUG 级别时有用）。

\paragraph{何时打开}

\begin{itemize}
  \item 站点数据筛选结果出乎意料 → 看 custom filter 的逐站日志
  \item \yamlkey{compare_tim_res} 等字段似乎没生效 → 看 processing 层的属性 fallback 警告
  \item 一般 evaluation/comparison 调试：\emph{用不上 \yamlkey{debug_mode}}，直接看 \file{run.log} 已经够
\end{itemize}

\subsection{日志去向}

\begin{itemize}
  \item 默认 stderr：可在终端实时看到
  \item \file{output/<run>/run.log}：完整记录（包括 debug）
\end{itemize}

错误诊断时常见做法：

\begin{minted}{bash}
# 把 stderr 与 stdout 都收集到文件
openbench run openbench.yaml 2>&1 | tee my_run.log

# 按级别过滤
grep "ERROR\|WARN" my_run.log
\end{minted}

\section{中断与续跑}

\subsection{Ctrl-C 中断后}

\openbench{} 在每个 task 完成时写 cache 与输出。Ctrl-C 不会破坏已完成 task 的结果，只会丢失正在进行的那个。

续跑：

\begin{minted}{bash}
openbench run openbench.yaml
# 增量缓存自动跳过已完成的 task；从中断处续上
\end{minted}

\subsection{异常崩溃后}

如果 run 因 OOM、磁盘满、HDF5 lock 等异常退出：

\begin{enumerate}
  \item 看日志末尾 traceback 定位失败的 task
  \item 修复根因（加内存、清磁盘、调小并行）
  \item 重跑相同命令；缓存自动跳过已完成的，从失败处继续
\end{enumerate}

\subsection{清缓存重跑（最后手段）}

如果增量缓存本身被污染（比如机器突然断电导致 partial write），强制清干净：

\begin{minted}{bash}
rm -rf output/<run>/scratch
openbench run openbench.yaml --cores 8
\end{minted}

或用 \yamlkey{project.force: true}。

\section{并行与性能}

\subsection{核数选择}

\begin{itemize}
  \item \yamlkey{project.num_cores: null}（默认）→ 自动 \texttt{os.cpu\_count()}
  \item 整数 → 严格用此核数
  \item CLI \texttt{--cores N} → 临时覆盖
\end{itemize}

\subsection{parallelism 现状}

当前版本（v3.0a1）的并行主要发生在两处：

\begin{itemize}
  \item DatasetProcessing 内部（站点处理、按年合并 NC）
  \item 部分 metric 计算用 dask
\end{itemize}

\noteinline{变量级并行（多个 task 并发执行）目前仍 \emph{保留接口未启用}。所以 \yamlkey{num_cores=16} 与 \yamlkey{num_cores=8} 在多变量评估时差距可能不显著。运维卷第 4 章会详细分析当前并行边界与未来路径。}

\subsection{HDF5 locking}

HDF5 ≥ 1.14 默认对只读 NC 文件加锁，多 worker 同时读同一参考文件时会卡死。\openbench{} 在 \cli{run} 启动时自动设置 \texttt{HDF5\_USE\_FILE\_LOCKING=FALSE} 解决这个问题。

\warninline{这个环境变量只在 \cli{openbench run} 进程内生效。如果你在 Jupyter 等环境里手动调用 \modname{openbench.runner.local.run\_evaluation}，需要自己 \texttt{os.environ["HDF5\_USE\_FILE\_LOCKING"] = "FALSE"} 或用 \texttt{xarray.open\_dataset(..., engine="h5netcdf", lock=False)}。}

\section{下一步}

下一章解读运行后的输出：

\begin{itemize}
  \item 输出目录布局每个文件夹的含义
  \item Metrics CSV 的列与解读
  \item Scores 的归一化方式
  \item Comparison 与 Ranking 输出
  \item Statistics 各模块产物
  \item HTML/PDF 报告结构
  \item 用 Python 直接二次分析结果
\end{itemize}
